\clearpage
\section*{《数学讲义丛书》编辑方针}\label{sec:publisher-policy}

\noindent\textit{Lecture Notes in Mathematics}\par
\noindent 主编：J.-M. Morel、B. Teissier

\begin{enumerate}
\item 《数学讲义丛书》旨在以迅速、非正式而高水平的方式，报道数学各领域及其应用的新进展. 欢迎分析建模和数值模拟新发展的数学著作.

书稿应当相当自足，内容完整. 因此，它们可以，而且经常不仅介绍作者自己的成果，也介绍他人的相关工作. 书稿可以以专题课程为基础. 此外，书稿还应提供充分的动机、例子与应用. 这使《数学讲义丛书》明显有别于通常十分简洁的期刊论文或技术报告. 那些原拟投期刊、只是篇幅太长而难以被多数期刊接受的文章，通常不具备这种“讲义”性质. 基于类似理由，本丛书通常不接收博士论文，但任教资格论文可能适合.

\item 除专著外，也欢迎由暑期学校或类似强化课程产生的多人合著书稿，前提是这些课程旨在向研究生初级或中级程度的听众介绍活跃的数学主题，并应提供参与者名单.

最终书稿不应只是课程笔记的汇编，而应要求主要讲师事先规划和协调. 全书结构应由主题内容决定，并在科学性导论中说明其动机与安排. 如有可能，编者应统一符号、参考文献、索引与结果的表述. 每篇内容应包含摘要，以及关联其他篇章的引言. 换言之，合著卷册需要的准备工作，远超过将活动期间提交的零散文章简单拼合.

\item 书稿应通过\url{https://www.editorialmanager.com/lnm}在线提交至Springer海德堡数学编辑部，或以电子方式提交给一位丛书编委. 作者应注意，不完整或距离最终稿不够接近的书稿，几乎总会导致更长的审稿时间和仍不明确的审稿建议，从而使最终稿还需再次审阅. 最低限度的材料应包括详细纲要，说明各章拟定内容，并附参考文献与若干样章. LNM审稿期间，不得同时向其他出版者投稿；平行投稿可能导致拒稿.

\item 一般而言，专著会送交至少两位外部审稿人评审. 只有依据完整书稿，才能作出最终出版决定；但旨在形成初步决定的审稿过程，可以以接近最终版本或尚未完成的书稿为基础.

多人合著卷册的分卷编者应安排按通常学术标准对各篇内容进行评审. 若能将审稿报告转交LNM编委会，将十分有帮助. 如果没有转交报告，或者统一性等其他问题仍不明确，丛书编委可能会咨询外部审稿人，对全卷作总体评估.

\item 书稿一般应使用英文提交. 最终稿应包含至少100页数学正文，并且始终包括：
\begin{itemize}
\item 目录；
\item 信息充分的引言，提供足够动机，并可酌情介绍一些历史背景，使不熟悉该主题的读者也能够理解；
\item 主题索引，通常这对读者很有帮助；
\item 为便于评审，书稿应以PDF文件形式提交.
\end{itemize}

\item 精心准备书稿，既能缩短制作周期，也能确保纸质版和在线版具有令人满意的外观. 接收后，作者将被要求准备最终LaTeX源文件（在线模板见\url{https://www.springer.com/gb/authors-editors/book-authors-editors/manuscriptpreparation/5636}）以及相应的PDF文件或压缩PS文件. LaTeX源文件是制作全书在线全文版本所必需的，现有LNM在线卷册见\url{http://link.springer.com/bookseries/304}. 《数学讲义丛书》一卷的技术制作约需12周. 如有必要，可通过\href{mailto:lnm@springer.com}{lnm@springer.com}索取进一步说明.

\item 作者共可获得30册免费样书，并免费访问SpringerLink上的本书，但不收取版税. 若直接从Springer订购供个人使用的Springer图书，可享受33.3\%折扣.

\item 出版承诺通过《出版协议》作出；多人合著卷册的撰稿作者须签署《同意出版书》. Springer-Verlag为每一卷登记版权. 作者可以自由地在后续出版物中再利用其LNM卷册所含材料：仅需书面或电子邮件简要请求正式许可.
\end{enumerate}

\noindent\textbf{联系地址}

\noindent Jean-Michel Morel教授，CMLA，法国卡尚高等师范学校\par
电子邮箱：\href{mailto:moreljeanmichel@gmail.com}{moreljeanmichel@gmail.com}

\noindent Bernard Teissier教授，几何与动力学组，法国巴黎朱西厄数学研究所——巴黎左岸\par
电子邮箱：\href{mailto:bernard.teissier@imj-prg.fr}{bernard.teissier@imj-prg.fr}

\noindent Springer：Ute McCrory，数学编辑部，德国海德堡\par
电子邮箱：\href{mailto:lnm@springer.com}{lnm@springer.com}
